\topic{vehicle-acceleration}{Vehicle acceleration: the rocket equation with gravity and drag}

Reading anchors: Hill and Peterson, second edition, Section 10.3,
printed pp.~472--477 / PDF pp.~484--489, Eqs.~(10.8)--(10.18)
\cite{hillpeterson}; Sutton and Biblarz, ninth edition, Section 4.1,
printed pp.~99--102, Eqs.~4--1 to 4--8, and Sections 4.2--4.3, from printed
p.~104, Eqs.~4--9 to 4--20 \cite{sutton}; Oates, third edition,
Section 3.3, ``Acceleration of a Rocket'', printed pp.~66--67,
Eqs.~(3.10)--(3.14) \cite{oates}.

The question is how fast a vehicle that is losing mass moves. The answer is
a momentum balance on a \emph{closed} system --- vehicle plus the mass it
expels --- and it settles the point that
Section~\ref{topic:force-one-form} left open: $F=d(Mu)/dt$ is wrong for a
rocket, and $M\,du/dt=F+\dots$, with $F$ the thrust, is right.

\subsection{Variables, assumptions and notation}
\label{sec:va-notation}
The model is flight over a non-rotating Earth.
\begin{description}\itemsep1pt
\item[$t$] time, s; the burn occupies $0\le t\le t_p$.
\item[$\mathbf u(t)$] vehicle velocity, m/s; $u=\lVert\mathbf u\rVert$ is
  the speed, $u_0=u(0)$, $u_p=u(t_p)$.
\item[$\hat e(t)$] unit vector along the flight path, $\mathbf u=u\hat e$;
  $\hat n$ is the in-plane unit normal, $\hat e$ turned by $+\pi/2$.
\item[$\hat e_T$] unit vector along the thrust (forward nozzle axis).
\item[$\theta,\ \psi$] angles of $\hat e$ and $\hat e_T$ above the local
  horizontal, rad.
\item[$M(t)$] vehicle mass, including the propellant still aboard, kg;
  $M(0)$ initial, $M_f=M(t_p)$ final (burnout).
\item[$m_p(t)$] propellant aboard; $M(t)=M_f+m_p(t)$
  (Eq.~\ref{Eq:MassBookkeeping}).
\item[$\dot m(t)$] propellant mass flow rate, $\dot m=-\frac{dM}{dt}\ge0$,
  kg/s.
\item[$u_{\rm exh},\ p_{\rm exh},\ A_e$] nozzle exit velocity relative to the
  vehicle (m/s), exit pressure (Pa) and exit area (m$^2$).
\item[$p_a$] ambient pressure, Pa.
\item[$F$] thrust, Eq.~\ref{Eq:ThrustEquation}, N.
\item[$u_e$] effective exhaust velocity, $u_e:=F/\dot m=I_sg_0$, m/s.
\item[$D,\ L$] drag and lift, N.
\item[$g$] local gravitational acceleration, m/s$^2$; $g_0$ the standard
  value.
\item[$\frac{M_f}{M(0)},\ \zeta$] mass ratio and propellant fraction
  $\frac{m_p(0)}{M(0)}=1-\frac{M_f}{M(0)}$, Eqs.~\ref{Eq:MassRatioDefinition},
  \ref{Eq:PropellantMassFractionDefinition}, \ref{Eq:ZetaPlusMR}.
\end{description}
Assumptions used throughout: the exhaust leaves uniformly along
$-\hat e_T$ with relative speed $u_{\rm exh}$; the vehicle does not rotate.
Each later specialisation states its additional assumptions.

The sources disagree on almost every symbol:
\begin{center}\small
\begin{tabular}{llll}
Ours & Sutton & Hill--Peterson & Oates\\ \hline
$M(t)$ & $m$ & $\mathcal M$ & $m$\\
$M(0)$, $M_f$ & $m_0$, $m_f$ & $\mathcal M_0$, $\mathcal M_b$
  & $m_0$, $m_0-m_p$\\
$\frac{M_f}{M(0)}$ & $\frac{m_f}{m_0}$ & $\frac1{\mathfrak R}$ & $\frac{m_0-m_p}{m_0}$\\
$u_e$ & $c$ & $u_{\rm eq}$ & $C$\\
$u_{\rm exh}$, $p_{\rm exh}$, $p_a$ & $v_2$, $p_2$, $p_3$
  & $u_e$, $p_e$, $p_a$ & $u_e$, $p_e$, $p_a$\\
$u$ & $u$ & $u$ & $v$\\
$F$, $D$ & $F$, $D$ & $\mathcal T$, $\mathfrak D$ & $T$, $D$\\
$\theta$ & $\theta$ & $\tfrac\pi2-\theta$ & $\tfrac\pi2-\theta$\\
$t_p$, $g_0$ & $t_p$, $g_0$ & $t_b$, $g_e$ & $\tau$, $g_0$\\
\end{tabular}
\end{center}
Four entries are traps.
\begin{itemize}\itemsep1pt
\item Hill and Peterson's $u_e$ is the actual exit velocity, our
  $u_{\rm exh}$. Our $u_e$ is the effective velocity, their $u_{\rm eq}$.
\item Hill and Peterson's mass ratio $\mathfrak R=\mathcal M_0/\mathcal M_b$
  and Oates's mass ratio $m_0/(m_0-m_p)$ are both $\frac{M(0)}{M_f}\ge1$.
\item Hill and Peterson and Oates measure $\theta$ from the local upward
  vertical, so their gravity term $-g\cos\theta_{\rm HP}$ is our
  $-g\sin\theta$. Hill and Peterson's text calls it ``the angle between
  $\mathbf u$ and $\mathbf g$'' (p.~473), but their Fig.~10.3 draws it
  between the tail axis, $-\mathbf u$, and $\mathbf g$; the figure is right.
\item Hill and Peterson's sounding-rocket example writes $u_e$ where its
  derivation needs $u_{\rm eq}$.
\end{itemize}

\subsection{The momentum balance, as Hill and Peterson argue it}
\label{sec:va-momentum}
Newton's second law applies to a fixed collection of matter. The vehicle
is not one: it loses mass. Hill and Peterson therefore follow a fixed
collection over a short interval. Fix $t$ and $dt>0$ (a small number here),
and take as the system $S$ everything aboard the vehicle at time $t$, of mass
$M$. During $[t,t+dt]$ the vehicle expels the parcel $dm=\dot m\,dt$, to
first order. At $t+dt$, $S$ consists of the vehicle, of mass $M-dm$ and speed
$u+du$, and the parcel. The parcel leaves at some instant in the interval,
with velocity $-u_{\rm exh}\hat e$ relative to the vehicle ($\psi=\theta$).
Its speed is therefore $u+\alpha\,du-u_{\rm exh}$ for some
$\alpha\in[0,1]$. Components are taken along the fixed vector $\hat e(t)$:
\[
\begin{aligned}
 P(t)&=Mu,\\
 P(t+dt)&=(M-dm)(u+du)+dm\,(u+\alpha\,du-u_{\rm exh}).
\end{aligned}
\]
Expanding the two products,
\[
\begin{aligned}
 (M-dm)(u+du)&=Mu+M\,du-dm\,u-dm\,du,\\
 dm\,(u+\alpha\,du-u_{\rm exh})&=dm\,u+\alpha\,dm\,du-dm\,u_{\rm exh} .
\end{aligned}
\]
Add them and subtract $P(t)$. The terms $\pm dm\,u$ cancel:
\begin{equation}\label{eq:va-momentum-change}
 P(t+dt)-P(t)=M\,du-dm\,u_{\rm exh}-(1-\alpha)\,dm\,du .
\end{equation}
The last term is second order and is dropped.

The external forces on $S$ along $\hat e$ are:
\begin{description}\itemsep1pt
\item[Pressure.] The ambient pressure is $p_a$ over the closed surface of
  $S$, except on the jet cross-section $A_e$, where the earlier exhaust
  presses with $p_{\rm exh}$. As in the derivation of
  Eq.~\ref{Eq:ThrustEquation}, a uniform $p_a$ over a closed surface
  integrates to zero. This leaves $(p_{\rm exh}-p_a)A_e$ along $+\hat e$.
\item[Drag.] $-D$, which collects the viscous traction and the departure of
  the external pressure from $p_a$.
\item[Gravity.] $-Mg\sin\theta$.
\end{description}
Equating the impulse over $dt$ to Eq.~\ref{eq:va-momentum-change},
substituting $dm=\dot m\,dt$, dividing by $dt$ and letting $dt\to0$,
\[
 M\frac{du}{dt}-\dot m\,u_{\rm exh}=(p_{\rm exh}-p_a)A_e-D-Mg\sin\theta .
\]
Moving $\dot m\,u_{\rm exh}$ to the right and using
Eq.~\ref{Eq:ThrustEquation},
\begin{equation}\label{eq:va-tangential}
\begin{aligned}
 M\frac{du}{dt}
 &=\dot m\,u_{\rm exh}+(p_{\rm exh}-p_a)A_e-D-Mg\sin\theta\\
 &=F-D-Mg\sin\theta .
\end{aligned}
\end{equation}

\paragraph{Sources side by side.} Hill and Peterson's
\[
 \mathcal M\,du-dm\,u_e
 =\bigl[(p_e-p_a)A_e-\mathfrak D-\mathcal Mg\cos\theta_{\rm HP}\bigr]dt
 \qquad\text{\cite[Eq.~(10.8)]{hillpeterson}}
\]
is Eq.~\ref{eq:va-tangential} times $dt$. They assign the parcel the speed
$u$ before and $u-u_e$ after, which differs from
Eq.~\ref{eq:va-momentum-change} only in the second-order $\alpha$ term.
Oates writes the same balance directly,
$T-D-gm\cos\theta=m\,dv/dt$ \cite[Eq.~(3.10)]{oates}. Sutton's Eq.~4--13,
$m\,du/dt=F\cos(\psi-\theta)-D-mg\sin\theta$, is Eq.~\ref{eq:va-tangential}
with $\psi=\theta$; his Eq.~4--1, $F=m\,du/dt$, is its gravity-free,
drag-free case. He calls it Newton's second law, but it is not $d(mu)/dt$.

\paragraph{Why $F=d(Mu)/dt$ fails.} By the product rule,
\[
 \frac{d}{dt}(Mu)=M\frac{du}{dt}+u\frac{dM}{dt}
 =M\frac{du}{dt}-\dot m\,u .
\]
Change to another inertial frame, $u\mapsto u+w$ with $w$ constant. Then
$M\,du/dt$ is unchanged, but $d(Mu)/dt$ changes by $-\dot m\,w$. Forces do
not depend on the inertial frame, so $F=d(Mu)/dt$ cannot hold when
$\dot m\neq0$. Equation~\ref{eq:va-tangential} contains only the relative
exhaust velocity, so it is frame-independent.

\subsection{The same balance, rigorously}
\label{sec:va-rigour}
\subsubsection*{What the increment argument proves}
Read literally, the display for $P(t+dt)$ has four problems.
\begin{enumerate}\itemsep1pt
\item $dt$, $dm$, $du$ are finite increments $h$, $\Delta m$, $\Delta u$,
  not differentials. The display is not an identity; it is a first-order
  expansion with an unstated remainder.
\item The parcel does not leave at one instant. Mass leaves continuously
  over $[t,t+h]$, so its momentum is an integral over emission times, and
  after leaving it is pushed by the jet and by gravity for a time at most
  $h$.
\item Momentum is a vector. For a vector-valued $\mathbf u(t)$ there is in
  general \emph{no} $\alpha\in[0,1]$ with
  $\mathbf u(\tau)=\mathbf u(t)+\alpha(\mathbf u(t+h)-\mathbf u(t))$: the
  mean value theorem fails for curves that do not lie on a line, for
  example $\mathbf u(t)=(\cos t,\sin t)$. What is true, and all that is
  needed, is $\lVert\mathbf u(\tau)-\mathbf u(t)\rVert\le
  h\sup\lVert\dot{\mathbf u}\rVert$.
\item Taking components along $\hat e(t)$ is legitimate only because
  $\hat e(t)$ is held fixed: pairing with a fixed vector is linear and
  commutes with $d/dh$.
\end{enumerate}
With $M$ and $\mathbf u$ of class $C^1$ and bounded forces, items 2 and 3
contribute terms of norm at most $Ch^2$. So the honest content of the
display is
\[
 \mathbf P_S(t+h)-\mathbf P_S(t)
 =M\,(\mathbf u(t+h)-\mathbf u(t))-\Delta m\,u_{\rm exh}\hat e_T
  +\mathbf R(h),
 \qquad \lVert\mathbf R(h)\rVert\le Ch^2,
\]
which says that $h\mapsto\mathbf P_S(t+h)$ is differentiable at $h=0$ with
derivative $M\dot{\mathbf u}-\dot m\,u_{\rm exh}\hat e_T$. That derivative is
all the argument needs, and all it can deliver. The derivation below obtains
it exactly, without increments.

\subsubsection*{Where $\mathbf P$ lives, and what $dt$ is}
\begin{description}\itemsep2pt
\item[Spacetime and time.] Galilean spacetime is a $4$-dimensional affine
  space $\mathbb A$ with translation space $W$ and a nonzero linear map
  $W\to\R$, the clock. The clock defines a closed, nowhere-vanishing
  $1$-form $dt$ on $\mathbb A$. Its kernel $V\subset W$, the spatial
  translations, carries the Euclidean inner product, and the level sets
  $t=\text{const}$ are Euclidean $3$-spaces $E$ with translation space $V$.
  On the time line $I\subset\R$, $dt$ is the coordinate $1$-form,
  $dt(\partial_t)=1$. So yes, $dt$ is a $1$-form. The ``small $dt$'' of the
  increment argument is the number $dt(h\,\partial_t)=h$.
\item[Velocity and frames.] The vehicle's world line is a curve
  $\gamma:I\to\mathbb A$ with $dt(\gamma')=1$. An inertial frame is a choice
  of $e_0\in W$ with $dt(e_0)=1$, and then $\mathbf u=\gamma'-e_0\in V$.
  Changing the frame, $e_0\mapsto e_0-\mathbf w$, adds $\mathbf w\in V$ to
  every velocity: the Galilean boost of the previous paragraph.
\item[Momentum.] Because $E$ is affine, every tangent space $T_xE$ is
  canonically $V$. Momenta of particles at \emph{different} points can
  therefore be added, which is the flatness used in
  Proposition~\ref{prop:impulse-momentum}. Formally: the translation group
  $(V,+)$ acts on the configuration space $Q=E^N$ of $N$ particles, its Lie
  algebra is $V$, and the momentum map of the induced action on the
  symplectic manifold $T^*Q$ is
  \[
   J:T^*Q\to V^*,\qquad
   \langle J(q,p),\boldsymbol\xi\rangle=\sum_a p_a(\boldsymbol\xi),
  \]
  each $p_a\in T^*_{q_a}E\cong V^*$. For $p_a=m_a\mathbf v_a^\flat$
  (Proposition~\ref{prop:momentum-covector}), $J=\mathbf P^\flat$ with
  $\mathbf P=\sum_am_a\mathbf v_a$. For a continuum,
  $\mathbf P=\int\rho\,\mathbf v\,\mathrm{vol}$. So the total momentum of a
  system is an element of $V^*$, identified with $V$ by the metric. Its time
  history is a smooth curve $\mathbf P:I\to V^*$, and its differential
  $d\mathbf P=\dot{\mathbf P}\,dt$ is a $V^*$-valued $1$-form on $I$. In a
  curved space there is no translation group, and no total momentum.
\item[Newton's law as forms.] For a closed system, $d\mathbf P=\mathbf F\,dt$
  in $\Omega^1(I;V^*)$, with $\mathbf F$ the external force. This is the
  forced Noether balance, Eq.~\ref{Eq:ForcedSymmetryBalance}, applied to every
  translation generator $\boldsymbol\xi\in V$ at once. Integrating over
  $[t_0,t_1]$ (Stokes on a $1$-manifold with boundary) gives the impulse
  identity.
\item[Products of differentials.] Hill and Peterson's Eq.~(10.8),
  $\mathcal M\,du-u_e\,dm=[\dots]\,dt$, is an exact identity between
  $1$-forms on $I$, with $dm:=\dot m\,dt=-d\mathcal M$. Evaluating it on
  $h\,\partial_t$ gives the increment reading to first order. A product such
  as $dm\,du$ is not a $1$-form. As a tensor it is $dm\otimes du$, quadratic
  in $h$, and its antisymmetric part $dm\wedge du$ vanishes identically
  because $\Lambda^2T^*I=0$ on a $1$-manifold. Dropping it is exactly what
  taking the differential means.
\end{description}

\subsubsection*{Derivation from the balance of momentum of a material body}
Let the matter be a continuum with density $\rho(t,x)$ and spatial velocity
field $\mathbf v(t,x)$, and let $\varphi_s$ be its motion, so that a material
region $\mathcal P$ occupies $\varphi_s(\mathcal P)\subset E$ at time $s$. The
axiom is Euler's balance of linear momentum: for every material region,
\begin{equation}\label{eq:va-euler-axiom}
 \frac{d}{ds}\int_{\varphi_s(\mathcal P)}\rho\,\mathbf v\,\mathrm{vol}
 =\oint_{\partial\varphi_s(\mathcal P)}\mathbf t\,dA
  +\int_{\varphi_s(\mathcal P)}\rho\,\mathbf g\,\mathrm{vol},
\end{equation}
with $\mathbf t$ the Cauchy traction. The integrands are $V$-valued
$3$-forms on $E$; they are integrated componentwise in the fixed space $V$.

\paragraph{Transport theorem.} Let $\Omega_s=\Phi_s(\Omega_t)$, where
$\Phi_s$ is the flow of a time-dependent vector field $\mathbf w$ with
$\Phi_t=\mathrm{id}$, and let $\omega_s=f(s,\cdot)\,\mathrm{vol}$ with $f$
$V$-valued. Pull back to the fixed region, differentiate, and use Cartan's
formula $\mathcal L_{\mathbf w}=d\,\iota_{\mathbf w}+\iota_{\mathbf w}d$,
with $d\omega=0$ in top degree, and Stokes:
\begin{equation}\label{eq:va-transport}
\begin{aligned}
 \frac{d}{ds}\Big|_{s=t}\int_{\Omega_s}\omega_s
 &=\frac{d}{ds}\Big|_{s=t}\int_{\Omega_t}\Phi_s^*\omega_s
 =\int_{\Omega_t}\Bigl(\frac{\partial\omega}{\partial s}
   +\mathcal L_{\mathbf w}\omega\Bigr)\\
 &=\int_{\Omega_t}\frac{\partial f}{\partial t}\,\mathrm{vol}
   +\oint_{\partial\Omega_t}f\,\iota_{\mathbf w}\mathrm{vol}.
\end{aligned}
\end{equation}
On $\partial\Omega_t$, $\iota_{\mathbf w}\mathrm{vol}=\langle\mathbf w,
\mathbf n\rangle\,dA$ with $\mathbf n$ the outward normal.

\paragraph{The vehicle as a moving region.} Let $B_s$ be the region occupied
by the vehicle and its internal contents. Without rotation it translates:
its boundary moves with the vehicle velocity $\mathbf u(s)$. Fix $t$, and
let $\mathcal P$ be the matter inside the vehicle at that instant,
$\varphi_t(\mathcal P)=B_t$. Apply Eq.~\ref{eq:va-transport} twice at
$s=t$: to $\varphi_s(\mathcal P)$, with $\mathbf w=\mathbf v$, and to $B_s$,
with $\mathbf w=\mathbf u$. The two regions coincide at $s=t$, so the volume
integrals agree, and subtracting,
\begin{equation}\label{eq:va-material-vs-region}
 \frac{d}{ds}\Big|_{s=t}\int_{\varphi_s(\mathcal P)}\rho f\,\mathrm{vol}
 =\frac{d}{ds}\Big|_{s=t}\int_{B_s}\rho f\,\mathrm{vol}
  +\oint_{\partial B_t}\rho f\,\langle\mathbf v-\mathbf u,\mathbf n\rangle\,dA .
\end{equation}
Matter crosses $\partial B_t$ only through the exit plane $A_e$, where
$\mathbf n=-\hat e_T$ and $\mathbf v-\mathbf u=-u_{\rm exh}\hat e_T$, so
$\langle\mathbf v-\mathbf u,\mathbf n\rangle=u_{\rm exh}$.

\emph{Mass}, $f=1$. The mass of $\mathcal P$ is constant, so
\[
 0=\frac{dM}{dt}+\int_{A_e}\rho\,u_{\rm exh}\,dA
 =\frac{dM}{dt}+\dot m,
 \qquad \dot m:=\rho_{\rm exh}u_{\rm exh}A_e .
\]

\emph{Momentum}, $f=\mathbf v$. Write the momentum inside $B_s$ as
$\int_{B_s}\rho\mathbf v\,\mathrm{vol}=M\mathbf u+\mathbf Q$, where
$\mathbf Q:=\int_{B_s}\rho(\mathbf v-\mathbf u)\,\mathrm{vol}$ is the momentum
of the contents relative to the vehicle. Assume $\mathbf Q$ constant:
propellant at rest in the tanks contributes nothing as it drains, and the
feed, chamber and nozzle flows are quasi-steady. Then
\[
\begin{aligned}
 \frac{d}{dt}\int_{\varphi_t(\mathcal P)}\rho\mathbf v\,\mathrm{vol}
 &=\frac{dM}{dt}\mathbf u+M\dot{\mathbf u}
   +\int_{A_e}\rho\,(\mathbf u-u_{\rm exh}\hat e_T)\,u_{\rm exh}\,dA\\
 &=-\dot m\,\mathbf u+M\dot{\mathbf u}
   +\dot m\,\mathbf u-\dot m\,u_{\rm exh}\hat e_T\\
 &=M\dot{\mathbf u}-\dot m\,u_{\rm exh}\hat e_T .
\end{aligned}
\]
The cancelling $\pm\dot m\,\mathbf u$ is Hill and Peterson's cancellation of
$\pm dm\,u$, and it is why the result is frame-independent.

\emph{Forces.} On $\partial B_t\setminus A_e$ split the traction into the
ambient part and the aerodynamic part, and define the aerodynamic resultant
\[
 \mathbf A:=\oint_{\partial B_t\setminus A_e}(\mathbf t+p_a\mathbf n)\,dA
 =-D\,\hat e+L\,\hat n .
\]
On $A_e$, $\mathbf t=-p_{\rm exh}\mathbf n$. Since $p_a$ is constant,
$\oint_{\partial B_t}p_a\mathbf n\,dA=0$ by Stokes, and
\[
 \oint_{\partial B_t}\mathbf t\,dA
 =-\int_{A_e}(p_{\rm exh}-p_a)\,\mathbf n\,dA+\mathbf A
 =(p_{\rm exh}-p_a)A_e\,\hat e_T+\mathbf A .
\]
With uniform $\mathbf g$, the body force is $M\mathbf g$.

\emph{Result.} Substituting into Eq.~\ref{eq:va-euler-axiom} and moving the
flux term to the right,
\begin{equation}\label{eq:va-vector}
 M\dot{\mathbf u}
 =\bigl[\dot m\,u_{\rm exh}+(p_{\rm exh}-p_a)A_e\bigr]\hat e_T
  +\mathbf A+M\mathbf g
 =F\,\hat e_T-D\,\hat e+L\,\hat n+M\mathbf g .
\end{equation}

\paragraph{Planar components.} With $\hat e=(\cos\theta,\sin\theta)$,
$\hat n=(-\sin\theta,\cos\theta)$ and $\mathbf g=(0,-g)$, differentiate
$\mathbf u=u\hat e$:
$\dot{\mathbf u}=\dot u\,\hat e+u\dot\theta\,\hat n$. The pairings are
$\langle\hat e_T,\hat e\rangle=\cos(\psi-\theta)$,
$\langle\hat e_T,\hat n\rangle=\sin(\psi-\theta)$,
$\langle\mathbf g,\hat e\rangle=-g\sin\theta$ and
$\langle\mathbf g,\hat n\rangle=-g\cos\theta$. Pair
Eq.~\ref{eq:va-vector} with $\hat e$ and with $\hat n$:
\begin{equation}\label{eq:va-planar}
\begin{aligned}
 M\frac{du}{dt}&=F\cos(\psi-\theta)-D-Mg\sin\theta,\\
 Mu\frac{d\theta}{dt}&=F\sin(\psi-\theta)+L-Mg\cos\theta .
\end{aligned}
\end{equation}
These are Sutton's Eqs.~4--13 and 4--14. With $\psi=\theta$ the first is
Eq.~\ref{eq:va-tangential}. The vector form gives both at once; the scalar
argument gives only the first. A rotating vehicle adds terms (jet damping,
Coriolis forces on the internal flow) because its boundary velocity is then
$\mathbf u+\boldsymbol\omega\times\mathbf r$.

\subsubsection*{Free space: conservation of the momentum map}
With no external force, the closed system of vehicle plus all exhaust has
constant total momentum. Let $\mathbf U(\tau,t)$ be the velocity at time $t$
of the exhaust shed at time $\tau$, so $\mathbf U(\tau,\tau)=\mathbf u(\tau)
-u_{\rm exh}\hat e_T$. Then
\[
 \mathbf P(t)=M(t)\mathbf u(t)+\int_0^t\dot m(\tau)\,\mathbf U(\tau,t)\,d\tau .
\]
Differentiate, using the Leibniz rule for the integral:
\[
 0=\frac{d\mathbf P}{dt}
 =-\dot m\,\mathbf u+M\dot{\mathbf u}
  +\dot m\,(\mathbf u-u_{\rm exh}\hat e_T)
  +\int_0^t\dot m(\tau)\frac{\partial\mathbf U}{\partial t}(\tau,t)\,d\tau .
\]
The last integral is the rate of change of momentum of the exhaust already
outside, which in vacuum is the push $-p_{\rm exh}A_e\hat e_T$ it receives
across the exit plane. Hence $M\dot{\mathbf u}=(\dot m\,u_{\rm exh}
+p_{\rm exh}A_e)\hat e_T=F\hat e_T$ with $p_a=0$, which is
Eq.~\ref{eq:va-vector} in free space. If the exhaust is taken to coast after
leaving, the integral is zero and the relative speed is the effective $u_e$.

\paragraph{A Lagrangian route, and why it needs care.} Ernest's AE121
Problem Set~2 solution writes $M(t)=M(0)-m(t)$, with $m(t)$ the mass
expelled, and the ``Lagrangian''
\[
 \mathcal L=\tfrac12Mu^2+\tfrac12\int_0^t\dot m\,d\tau\,(u-u_e)^2 ,
\]
the second term the kinetic energy of the exhaust, $u_e$ the exhaust speed
relative to the vehicle. The Euler--Lagrange equation then gives
$-\dot mu_i+M\dot u_i+\dot mu_i-\dot mu_e=0$, that is
\[
 \boxed{\dot u_i=\frac{\dot mu_e}{M}} .
\]
The result is right, but two steps cancel each other. First, every past
parcel is given the current speed $u(t)-u_e$; the parcel shed at $\tau$
moves at $u(\tau)-u_e$. Second, in
\[
 \frac{d}{dt}\int_0^t\dot m(\tau)\bigl(u(t)-u_e\bigr)d\tau
 =\dot m(t)\bigl(u(t)-u_e\bigr)+\Bigl(\int_0^t\dot m\,d\tau\Bigr)\dot u ,
\]
the second term was dropped. Keeping it, without the first fix, gives
$M(0)\dot u=\dot mu_e$, a constant-mass equation. The vehicle alone has no
Lagrangian, because it is not a closed system. The correct object is the
momentum map of the closed system above, with coasting exhaust:
\[
 \frac{d}{dt}\Bigl[Mu+\int_0^t\dot m(\tau)\bigl(u(\tau)-u_e\bigr)d\tau\Bigr]
 =-\dot mu+M\dot u+\dot m\,(u-u_e)=M\dot u-\dot mu_e=0 .
\]

\subsection{Integrating: the ideal increment and the losses}
\label{sec:va-integration}
Divide Eq.~\ref{eq:va-tangential} by $M$ and write
$F=u_e\dot m=-u_e\frac{dM}{dt}$:
\begin{equation}\label{eq:va-acceleration}
 \frac{du}{dt}=-\frac{u_e}{M}\frac{dM}{dt}-\frac{D}{M}-g\sin\theta .
\end{equation}
This is Hill and Peterson's Eq.~(10.9), Oates's Eq.~(3.11), and Sutton's
Eq.~4--15 with $\psi=\theta$.

\paragraph{Gravity-free, drag-free, constant flow.} Sutton's Section~4.1
route: with $\dot m=m_p/t_p$ constant, $m(t)=\frac{m_p}{t_p}t$ expelled and
\[
 M(t)=M(0)-\frac{m_p}{t_p}t
 =M(0)\Bigl(1-\frac{m_p}{M(0)}\frac{t}{t_p}\Bigr)
 =M(0)\Bigl(1-\zeta\frac{t}{t_p}\Bigr),
\]
Sutton's Eqs.~4--2 and 4--3. Then $F=\dot mu_e=M\,du/dt$ gives
$\frac{\dot m\,dt}{M}=\frac{du}{u_e}$, and
\[
\begin{aligned}
 \frac{\Delta u}{u_e}
 &=\int_0^{t_p}\frac{\frac{m_p}{t_p}\,dt}{M(0)-\frac{m_p}{t_p}t}
 =\Bigl[-\ln\Bigl(M(0)-\frac{m_p}{t_p}t\Bigr)\Bigr]_0^{t_p}\\
 &=\ln M(0)-\ln\bigl(M(0)-m_p\bigr)
 =\ln\frac{M(0)}{M(0)-m_p}=\ln\frac{M(0)}{M_f}.
\end{aligned}
\]
Thus
\begin{equation}\label{eq:va-ideal}
\begin{aligned}
 \Delta u&=u_e\ln\frac{M(0)}{M_f},\\
 e^{\Delta u/u_e}&=\frac{M(0)}{M_f},\\
 \frac{M_f}{M(0)}&=e^{-\Delta u/u_e}.
\end{aligned}
\end{equation}
Sutton's Eqs.~4--5 to 4--7, Hill and Peterson's Eqs.~(10.10)--(10.11)
($\Delta u=u_{\rm eq}\ln\mathfrak R$), and Oates's Eq.~(3.14). Solving for
the propellant, $m_p=M(0)-M_f$,
\begin{equation}\label{eq:va-propellant}
 m_p=M_f\bigl(e^{\Delta u/u_e}-1\bigr)=M(0)\bigl(1-e^{-\Delta u/u_e}\bigr),
\end{equation}
Sutton's Eq.~4--8. By Eq.~\ref{Eq:ZetaPlusMR}, $\zeta=1-\frac{M_f}{M(0)}$; for
example $M(0)/M_f=5$ gives $\zeta=1-\frac15=0.8$. Constant flow was a
convenience: $\int\dot m\,dt/M=-\int dM/M$ for any flow history.

\paragraph{Payload and dead weight.} Oates splits the final mass into the
payload $m_L$ and the dead weight $m_d$ (structure, engines, unused
propellant), $M_f=m_L+m_d$ \cite[p.~67]{oates}. With the payload ratio
$\lambda:=m_L/M(0)$ and the dead-weight ratio $\delta:=m_d/M(0)$,
\[
 \frac{M_f}{M(0)}=\frac{m_L+m_d}{M(0)}=\lambda+\delta,
 \qquad
 \frac{M(0)}{M(0)-m_p}=\frac1{\lambda+\delta}.
\]

\paragraph{Keeping the pressure thrust explicit.} Oates's route: with
$p_a=0$ and $D=0$, $g=0$, and $\dot m=m_p/t_p$ constant, integrate the two
parts of $F/M$ separately:
\[
\begin{aligned}
 \Delta u
 &=-u_{\rm exh}\ln\frac{M(0)-m_p}{M(0)}
  +\Bigl[\frac{(p_{\rm exh}-p_a)A_e}{-m_p/t_p}
   \ln\Bigl(M(0)-\frac{m_p}{t_p}t\Bigr)\Bigr]_0^{t_p}\\
 &=-u_{\rm exh}\ln\frac{M(0)-m_p}{M(0)}
  +\frac{(p_{\rm exh}-p_a)A_e}{-\dot m}\ln\frac{M(0)-m_p}{M(0)}\\
 &=\Bigl[u_{\rm exh}+\frac{(p_{\rm exh}-p_a)A_e}{\dot m}\Bigr]
   \ln\frac{M(0)}{M(0)-m_p}
 =u_e\ln\frac{M(0)}{M_f} .
\end{aligned}
\]
The bracket is $u_e$ by Eq.~\ref{Eq:ThrustEquation}; the pressure term has
the same logarithm only because $(p_{\rm exh}-p_a)A_e/\dot m$ was constant.

\paragraph{With gravity and drag.} Integrate Eq.~\ref{eq:va-acceleration}
over the burn and define the two losses
\begin{equation}\label{eq:va-losses}
 \Delta u_g:=\int_0^{t_p}g\sin\theta\,dt,\qquad
 \Delta u_D:=\int_0^{t_p}\frac{D}{M}\,dt .
\end{equation}
With $u_e$ constant the first term is
$-u_e\int_{M(0)}^{M_f}dM/M=u_e\ln\frac{M(0)}{M_f}$, so
\begin{equation}\label{eq:va-budget}
 u_p-u_0=u_e\ln\frac{M(0)}{M_f}-\Delta u_g-\Delta u_D .
\end{equation}
This is Oates's Eq.~(3.13) and Sutton's Eq.~4--19.
\begin{itemize}\itemsep1pt
\item $D=0$, $g$ constant: $\Delta u_g=g\,t_p\,\overline{\sin\theta}$, the
  time average over the burn. This is Hill and Peterson's Eq.~(10.13),
  $\Delta u=u_{\rm eq}\ln\mathfrak R-g\,\overline{\cos\theta}\,t_b$, in their
  angle. For variable $g$, $\Delta u_g=t_p\,\overline{g\sin\theta}$, the
  average of the product.
\item Vertical flight, constant $\dot m$: with
  $D=C_D\tfrac12\rho u^2A$ and $M=M(0)(1-\zeta t/t_p)$,
  \[
   \Delta u_D=\int_0^{t_p}
   \frac{C_D\tfrac12\rho u^2A}{M(0)(1-\zeta t/t_p)}\,dt
   =\frac{C_DA}{M(0)}B,
   \quad
   B:=\int_0^{t_p}\frac{\tfrac12\rho u^2}{1-\zeta t/t_p}\,dt ,
  \]
  the drag term of Sutton's Eq.~4--19. Taking $C_D$ outside the integral
  assumes it is constant, but it depends on Mach number.
\end{itemize}
In ascent $\sin\theta>0$ and $D\ge0$, so both losses are nonnegative and
$u_p-u_0\le u_e\ln\frac{M(0)}{M_f}$.

\paragraph{When $u_e$ is not constant.} The first term is
$\int_0^{t_p}u_e\,\dot m/M\,dt$. Sutton's $\bar c$ in Eq.~4--19 is the
average that makes $\bar u_e\ln\frac{M(0)}{M_f}$ equal to it:
\[
 \bar u_e:=\frac{\int_0^{t_p}u_e\,\frac{\dot m}{M}\,dt}
          {\int_0^{t_p}\frac{\dot m}{M}\,dt},
 \qquad
 \int_0^{t_p}\frac{\dot m}{M}\,dt=\ln\frac{M(0)}{M_f} .
\]
The weight is $\dot m/M$, not the $\dot m$ of
Proposition~\ref{prop:weighted-mean}: propellant burned late, when the
vehicle is light, counts for more. In the atmosphere $u_e=u_{\rm exh}
+(p_{\rm exh}-p_a)A_e/\dot m$ rises as $p_a$ falls with altitude.

\subsection{Exercises: gravity-free flight}
\begin{exercise}[Sutton and Biblarz 9e, Chapter 4, Problem 1 \cite{sutton}]
For a vehicle in gravitationless space, determine the mass ratio necessary
to boost the vehicle velocity by (a) 1600 m/sec and (b) 3400 m/sec; the
effective exhaust velocity is 2000 m/sec. If the initial total vehicle mass
is 4000 kg, what are the corresponding propellant masses?
\emph{Answer:} (a) 2204 kg.
\end{exercise}
\begin{proof}[Solution]
By Eq.~\ref{eq:va-ideal}, $\frac{M_f}{M(0)}=e^{-\Delta u/u_e}$, and by
Eq.~\ref{eq:va-propellant}, $m_p=M(0)\bigl(1-\frac{M_f}{M(0)}\bigr)$.
\[
\begin{aligned}
 \text{(a)}&\quad \frac{M_f}{M(0)}=e^{-1600/2000}=e^{-0.8}=0.449329,\\
 &\quad m_p=4000\times0.550671=2202.7\ \mathrm{kg};\\
 \text{(b)}&\quad \frac{M_f}{M(0)}=e^{-3400/2000}=e^{-1.7}=0.182684,\\
 &\quad m_p=4000\times0.817316=3269.3\ \mathrm{kg}.
\end{aligned}
\]
The printed 2204 kg follows from rounding $\frac{M_f}{M(0)}$ to $0.449$.
\end{proof}

\begin{exercise}[Sutton and Biblarz 9e, Chapter 4, Problem 4 \cite{sutton}]
A research space vehicle in gravity-free and drag-free space launches a
smaller spacecraft into a meteor shower region. The 2-kg sensitive
instrument package of this spacecraft (25 kg total mass) limits the maximum
acceleration to no more than 50 m/sec$^2$. It is launched by a solid
propellant rocket motor ($I_s=260$ sec and $\zeta=0.88$). Assume
instantaneous start and stop of the rocket motor.
(a) Determine the minimum allowable burn time, assuming steady constant
propellant mass flow;
(b) Determine the maximum velocity relative to the launch vehicle.
(c) Solve for (a) and (b) when half of the total impulse is delivered at
the previous propellant mass flow rate, with the other half at 20\% of this
mass flow rate.
\end{exercise}
\paragraph{Solution.} Take $\zeta$ for the whole 25 kg spacecraft. With
constant thrust $F=I_sg_0\dot m=I_sg_0\zeta M(0)/t_p$, the acceleration
$a=F/M(t)$ grows as the mass falls, so its maximum is at burnout
(Eq.~\ref{Eq:BurnoutAcceleration}):
\[
 a_f=\frac{F}{M_f}=\frac{I_sg_0\zeta}{t_p}\,\frac{M(0)}{M_f}
 =\frac{I_sg_0}{t_p}\,\frac{1}{M(0)/m_p-1}.
\]
(a) $a_f\le50$ is a \emph{lower} bound on the burn time, since a longer burn
lowers $a_f$:
\[
 t_p\ge\frac{I_sg_0}{50}\,\frac{1}{1/0.88-1}
 =\frac{260\times9.80665}{50}\times7.3333=373.96\ \mathrm{s}
\]
(373.71 s with $g_0=9.8$).
(b) By Eq.~\ref{eq:va-ideal},
$\Delta u=-I_sg_0\ln(1-\zeta)=2549.73\times\ln(1/0.12)=5406.1\ \mathrm{m/s}$
(5402.4 m/s with $g_0=9.8$). Part (c) is not yet worked.

\subsection{Forces in the atmosphere}
\label{sec:va-atmosphere}
\paragraph{Lift and drag.} With $A$ a reference (frontal) area and $\rho$ the
local air density,
\begin{equation}\label{eq:va-lift-drag}
 L=C_L\tfrac12\rho Au^2,\qquad D=C_D\tfrac12\rho Au^2 ,
\end{equation}
Sutton's Eqs.~4--10 and 4--11 and Hill and Peterson's Eq.~(10.14). The
coefficients depend on Mach number and on the angle of attack. Hill and
Peterson fit the density by $\rho(h)=a\exp(-bh^{1.15})$ with
$a=1.2\ \mathrm{kg/m^3}$, $b=2.9\times10^{-5}$ for $h$ in metres; it falls
to 1\% of its sea-level value by 30 km. Above 300 km the density varies by
up to a factor of two with solar activity and between day and night, the
main uncertainty in orbital drag (Sutton Section~4.2).

\paragraph{Gravity.} Neglect the Earth's oblateness and local variations.
Newton's law of gravitation for a mass $M$ at distance $r$ from the centre
of an Earth of mass $M_e$ and radius $R_0$, and its value at the surface,
\[
 F=\frac{GM_eM}{r^2},\qquad Mg_0=\frac{GM_eM}{R_0^2},
\]
divide to give, with $r=R_0+h$,
\begin{equation}\label{eq:va-gravity}
 g=g_0\frac{R_0^2}{r^2}=g_0\Bigl(\frac{R_0}{R_0+h}\Bigr)^2,
\end{equation}
Sutton's Eq.~4--12 and Hill and Peterson's Eq.~(10.12). At 100 miles,
$g=0.95g_0$.

\paragraph{Wingless vehicles.} With $L=0$, $\psi=\theta$ and constant
$\dot m$, Eqs.~\ref{eq:va-planar} divided by $M=M(0)(1-\zeta t/t_p)$ become
\begin{equation}\label{eq:va-wingless}
\begin{aligned}
 \frac{du}{dt}
 &=\frac{u_e\zeta/t_p}{1-\zeta t/t_p}-g\sin\theta
  -\frac{C_D\tfrac12\rho u^2A/M(0)}{1-\zeta t/t_p},\\
 u\frac{d\theta}{dt}&=-g\cos\theta ,
\end{aligned}
\end{equation}
Sutton's Eq.~4--17 and the reduced Eq.~4--16. The thrust term is
$F/M=(u_e\zeta M(0)/t_p)/M$, with $\zeta/t_p$ and not $\zeta t/t_p$ in the
numerator. Vertical flight, $\theta=\pi/2$, gives Sutton's Eq.~4--18.

\subsection{Turning, and why \texorpdfstring{$u_e\ln\frac{M(0)}{M_f}$}{u_e ln(M(0)/M_f)} is the budget}
\label{sec:va-budget}
The two lines of Eq.~\ref{eq:va-wingless} are a gravity turn: gravity alone
bends the path. Vertical flight, $\theta=\pi/2$, makes the second right side
vanish, so a gravity turn must start at some $\theta(0)<\pi/2$; that is Hill
and Peterson's initial angle $\theta_1$ from the vertical, their Fig.~10.5.

Hill and Peterson integrate the trajectory by stepping the velocity vector.
Over each $\delta t$ they hold $\mathfrak D/\mathcal M$, $g$ and $\theta$
fixed and add three increments, from thrust, drag and gravity. That is the
explicit Euler method applied to Eq.~\ref{eq:va-vector}. Its global error is
$O(\delta t)$, which is their remark that accuracy depends on $\delta t$
relative to $t_b$. A Runge--Kutta step (the Runge--Kutta topic of the master)
is the natural replacement.

The thrust increment over $\delta t$ is
$\delta\mathbf u_F=(F/M)\,\hat e_T\,\delta t$. To first order its magnitude
is $\lVert\delta\mathbf u_F\rVert=(u_e\dot m/M)\,\delta t=-u_e\,\delta M/M$.
Integrating over the burn,
\begin{equation}\label{eq:va-thrust-magnitude}
 \int_0^{t_p}\frac{F}{M}\,dt
 =-u_e\int_{M(0)}^{M_f}\frac{dM}{M}
 =u_e\ln\frac{M(0)}{M_f} ,
\end{equation}
for \emph{every} thrust-direction history. Their
$\sum\lvert\delta\mathbf u_{\mathcal T}\rvert=u_{\rm eq}\ln\mathfrak R$
(p.~476) is this equation. By the argument of
Proposition~\ref{prop:scalar-bound}, applied to $F\hat e_T/M$ in place of $f$,
\[
 \Bigl\lVert\int_0^{t_p}\frac{F}{M}\,\hat e_T\,dt\Bigr\rVert
 \le\int_0^{t_p}\frac{F}{M}\,dt
 =u_e\ln\frac{M(0)}{M_f} ,
\]
with equality exactly when $\hat e_T$ is constant whenever $F>0$. So
$u_e\ln\frac{M(0)}{M_f}$ is a scalar budget. A trajectory spends it on
velocity change, gravity loss, drag loss and, when thrust and path are not
aligned, the steering loss
$\int_0^{t_p}\frac{F}{M}\bigl(1-\cos(\psi-\theta)\bigr)dt$ of
Eq.~\ref{eq:va-planar}. That is the sense in which Hill and Peterson say
that Eq.~(10.11) ``characterize[s] the propulsion requirements for any
mission'': the trajectory computation supplies the losses.

\subsection{The vertical sounding rocket}
\label{sec:va-sounding}
Hill and Peterson, pp.~476--477, Eqs.~(10.15)--(10.18), and Sutton's
Eq.~4--20. Additional assumptions:
\begin{itemize}\itemsep1pt
\item vertical flight, $\theta=\pi/2$; no drag, $D=0$;
\item constant gravity, $g=g_0$; constant $u_e$;
\item constant $\dot m=\zeta M(0)/t_p$, so $M(t)=M(0)(1-\zeta t/t_p)$;
\item $u(0)=0$, $h(0)=0$, and the vehicle lifts off at $t=0$.
\end{itemize}
Lift-off at $t=0$ needs $F>M(0)g_0$. Since $F=u_e\dot m=u_e\zeta M(0)/t_p$,
this is
\begin{equation}\label{eq:va-liftoff}
 \frac{u_e\,\zeta}{g_0\,t_p}>1 .
\end{equation}
Without it the vehicle waits on the pad, and the formulas below do not
apply.

\paragraph{Velocity.} Integrate Eq.~\ref{eq:va-acceleration} from $0$ to
$t$, with $\theta=\pi/2$ and $D=0$:
\begin{equation}\label{eq:va-sounding-u}
 u(t)=u_e\ln\frac{M(0)}{M(t)}-g_0t
     =-u_e\ln\Bigl(1-\zeta\frac{t}{t_p}\Bigr)-g_0t .
\end{equation}
This is Hill and Peterson's Eq.~(10.15). Their
$1-(1-1/\mathfrak R)\,t/t_b$ is our $1-\zeta t/t_p$, since
$1-1/\mathfrak R=1-\frac{M_f}{M(0)}=\zeta$. At burnout,
\begin{equation}\label{eq:va-sounding-up}
 u_p=u_e\ln\frac{M(0)}{M_f}-g_0t_p ,
\end{equation}
which is Sutton's Eq.~4--20 with constant $g$.

\paragraph{Burnout altitude.} $h_p=\int_0^{t_p}u\,dt$. The gravity part
gives $-\tfrac12g_0t_p^2$. For the thrust part, substitute
$x=\zeta t/t_p$, so $dt=(t_p/\zeta)\,dx$:
\[
 \int_0^{t_p}-\ln\Bigl(1-\zeta\frac{t}{t_p}\Bigr)dt
 =\frac{t_p}{\zeta}\int_0^{\zeta}-\ln(1-x)\,dx .
\]
An antiderivative follows from the product rule,
\[
 \frac{d}{dx}\bigl[(1-x)\ln(1-x)+x\bigr]
 =-\ln(1-x)-1+1=-\ln(1-x),
\]
so the integral is $(1-\zeta)\ln(1-\zeta)+\zeta$. With $1-\zeta=\frac{M_f}{M(0)}$,
\[
 \frac{t_p}{\zeta}\bigl(\frac{M_f}{M(0)}\ln \frac{M_f}{M(0)}+\zeta\bigr)
 =t_p\Bigl(1+\frac{M_f}{M(0)-M_f}\ln\frac{M_f}{M(0)}\Bigr).
\]
Hence
\begin{equation}\label{eq:va-sounding-hp}
 h_p=u_e\,t_p\Bigl(1+\frac{M_f}{M(0)-M_f}\ln\frac{M_f}{M(0)}\Bigr)
     -\tfrac12g_0t_p^2 .
\end{equation}
To compare with Hill and Peterson's Eq.~(10.16), put
$\frac{M_f}{M(0)}=1/\mathfrak R$:
\[
 \frac{M_f}{M(0)-M_f}\ln\frac{M_f}{M(0)}
 =\frac{-\mathfrak R^{-1}\ln\mathfrak R}{(\mathfrak R-1)/\mathfrak R}
 =-\frac{\ln\mathfrak R}{\mathfrak R-1},
\]
and Eq.~\ref{eq:va-sounding-hp} becomes their
$h_b=-u_et_b\frac{\ln\mathfrak R}{\mathfrak R-1}+u_et_b-\tfrac12g_et_b^2$.
The lower limit of their $h_b=\int u\,dt$ is $0$.

\paragraph{Coast to apogee.} After burnout, $du/dt=-g_0$. Multiply by
$u=dh/dt$:
\[
 \frac{d}{dt}\Bigl(\tfrac12u^2\Bigr)=-g_0\frac{dh}{dt}
 \quad\Longrightarrow\quad
 \tfrac12u_p^2+g_0h_p=g_0h_{\max},
\]
since $u=0$ at apogee. This is Hill and Peterson's Eq.~(10.17),
$h_{\max}=h_p+u_p^2/(2g_0)$, and it requires $u_p>0$. Square
Eq.~\ref{eq:va-sounding-up}:
\[
 \frac{u_p^2}{2g_0}
 =\frac{u_e^2}{2g_0}\ln^2\frac{M(0)}{M_f}
  -u_e\,t_p\ln\frac{M(0)}{M_f}+\tfrac12g_0t_p^2 .
\]
Add Eq.~\ref{eq:va-sounding-hp}. The $\tfrac12g_0t_p^2$ terms cancel:
\[
\begin{aligned}
 h_{\max}&=\frac{u_e^2}{2g_0}\ln^2\frac{M(0)}{M_f}
  -u_e\,t_p\ln\frac{M(0)}{M_f}\\
 &\quad+u_e\,t_p
  +u_e\,t_p\frac{M_f}{M(0)-M_f}\ln\frac{M_f}{M(0)} .
\end{aligned}
\]
Write $\frac{M_f}{M(0)}\ln \frac{M_f}{M(0)}=-\frac{M_f}{M(0)}\ln\frac{M(0)}{M_f}$ and collect
the two logarithms:
\[
 \ln\frac{M(0)}{M_f}\Bigl(1+\frac{M_f}{M(0)-M_f}\Bigr)
 =\frac{M(0)}{M(0)-M_f}\ln\frac{M(0)}{M_f} .
\]
Therefore
\begin{equation}\label{eq:va-sounding-hmax}
\begin{aligned}
 h_{\max}&=\frac{u_e^2}{2g_0}\ln^2\frac{M(0)}{M_f}\\
 &\quad-u_e\,t_p\Bigl(\frac{M(0)}{M(0)-M_f}\ln\frac{M(0)}{M_f}-1\Bigr).
\end{aligned}
\end{equation}
With $\frac{M_f}{M(0)}=1/\mathfrak R$,
$\frac{M(0)}{M(0)-M_f}\ln\frac{M(0)}{M_f}=\mathfrak R\ln\mathfrak R/(\mathfrak R-1)$,
which is Hill and Peterson's Eq.~(10.18).

\paragraph{Burning time.} Differentiate
Eq.~\ref{eq:va-sounding-hmax} with $u_e$, $\frac{M_f}{M(0)}$ and $g_0$ held fixed:
\[
 \frac{\partial h_{\max}}{\partial t_p}
 =-u_e\Bigl(\frac{M(0)}{M(0)-M_f}\ln\frac{M(0)}{M_f}-1\Bigr).
\]
The bracket is positive. Let $y=\frac{M(0)}{M_f}>1$, so that
$\frac{M(0)}{M(0)-M_f}\ln\frac{M(0)}{M_f}=y\ln y/(y-1)$. Define
$\varphi(y):=y\ln y-(y-1)$. Then $\varphi(1)=0$ and
$\varphi'(y)=\ln y+1-1=\ln y>0$ for $y>1$, so $\varphi(y)>0$, that is,
$y\ln y/(y-1)>1$. A shorter burn therefore always gains altitude: it spends
less propellant lifting propellant (Hill and Peterson's ``Burning Time''
paragraph; Oates's gravity-loss remark). The first term of
Eq.~\ref{eq:va-sounding-hmax} is the impulsive limit $t_p\to0$, which cannot
be reached. Eq.~\ref{eq:va-liftoff} bounds $t_p$ from above; the largest
tolerable acceleration, reached at burnout, bounds it from below:
\[
\begin{aligned}
 a_f&=\frac{F}{M_f}-g_0
 =\frac{u_e\,\zeta}{t_p}\,\frac{M(0)}{M_f}-g_0\le a_{\max}\\
 &\iff t_p\ge\frac{u_e\,\zeta}{a_{\max}+g_0}\,\frac{M(0)}{M_f}.
\end{aligned}
\]

\begin{exercise}[Sutton and Biblarz 9e, Chapter 4, Problem 2 \cite{sutton}]
Determine the burnout velocity and burnout altitude for a dragless
projectile with the following parameters for a simplified vertical
trajectory: $\overline{c}=2209$ m/sec; $m_p/m_0=0.57$; $t_p=5.0$ sec; and
$u_0=0$; $h_0=0$. Select a relatively small diameter missile with $L/D$ of
10 and an average vehicle density of 1200 kg/m$^3$.
\end{exercise}
\begin{proof}[Solution]
In our notation $u_e=2209$ m/s, $\zeta=0.57$, $\frac{M_f}{M(0)}=0.43$, and we take
$g_0=9.8$ m/s$^2$; Eq.~\ref{eq:va-liftoff} holds with
$u_e\zeta/(g_0t_p)=25.7$. Equations~\ref{eq:va-sounding-up} and
\ref{eq:va-sounding-hp} give
\[
\begin{aligned}
 u_p&=2209\ln\frac1{0.43}-9.8\times5.0=1815.33\ \mathrm{m/s},\\
 h_p&=2209\times5.0\Bigl(1+\frac{0.43\ln0.43}{0.57}\Bigr)
      -\tfrac12\times9.8\times5.0^2=3890.38\ \mathrm{m}.
\end{aligned}
\]
A direct numerical integration of Eq.~\ref{eq:va-wingless} in SymPy gives
$u_p=1815.32988528061$ m/s and $h_p=3890.37850288891$ m, agreeing to every
printed digit. The length-to-diameter ratio and vehicle density are needed
only for the drag of Problem~3.
\end{proof}

\subsection{A ballistic trajectory with drag}
\label{sec:va-ballistic}
Ernest's AE121 Problem Set~2, Problem~4, is the unpowered case of
Eq.~\ref{eq:va-vector}: $F=0$, $L=0$, constant mass $M$, drag opposite the
velocity, and an exponential atmosphere $\rho=\rho_0e^{-y/H}$ with scale
height $H$. Since $D\hat e=C_D\tfrac12\rho Au^2\,\mathbf u/u$, in Cartesian
components $\mathbf u=(u_x,u_y)$,
\[
\begin{aligned}
 \dot u_x&=-k\,e^{-y/H}\bigl(u_x^2+u_y^2\bigr)^{1/2}u_x,\\
 \dot u_y&=-g-k\,e^{-y/H}\bigl(u_x^2+u_y^2\bigr)^{1/2}u_y,\\
 \dot x&=u_x,\qquad \dot y=u_y,
 \qquad k:=\frac{\rho_0C_DA}{2M}.
\end{aligned}
\]
The data are a 15 cm iron cannonball of density
$7.8\times10^3\ \mathrm{kg/m^3}$ ($M=13.78$ kg), $C_D=0.1$,
$\rho_0=1.225\ \mathrm{kg/m^3}$, $H=10$ km, $g=9.8\ \mathrm{m/s^2}$, and launch
at $300$ m/s and $50^\circ$, giving $k=7.853\times10^{-5}\ \mathrm{m^{-1}}$.
Without drag the range is $u_0^2\sin(2\theta_0)/g=9044$ m. Integrated with
an adaptive Dormand--Prince method to ground impact,
\begin{center}\small
\begin{tabular}{lccc}
Case & Range (m) & Flight time (s) & Apogee (m)\\ \hline
(a) $\rho=\rho_0e^{-y/H}$ & 6108 & 41.58 & 2124\\
(b) $\rho=\rho_0$ & 5887 & 41.23 & 2091\\
(c) no drag & 9044 & 46.90 & 2694\\
\end{tabular}
\end{center}
Drag costs almost 3 km of range. The thinning of the air with height is
worth 221 m, which matters for precise impact prediction.
